\begin{table}[H]
\centering
\caption{Standardized Effect Sizes for Main Outcomes}
\label{tab:sde}
\begin{adjustbox}{max width=\textwidth}
\begin{threeparttable}
\begin{tabular}{lcccccc}
\toprule
Outcome & $\hat{\beta}$ & SE & SD($Y$) & SDE & SE(SDE) & Classification \\
\midrule
\multicolumn{7}{l}{\textit{Panel A: Pooled}} \\
Employment (all industries) & -0.0023 & 0.0005 & 1.971 & -0.0065 & 0.0015 & Small negative \\
Employment (food services) & -0.0027 & 0.0006 & 1.854 & -0.0080 & 0.0017 & Small negative \\
Employment (retail trade) & -0.0003 & 0.0005 & 1.754 & -0.0011 & 0.0017 & Null \\
Employment (healthcare) & -0.0022 & 0.0008 & 1.852 & -0.0065 & 0.0024 & Small negative \\
\midrule
\multicolumn{7}{l}{\textit{Panel B: Heterogeneous}} \\
Employment (high-poverty counties) & -0.0040 & 0.0007 & 1.890 & -0.0101 & 0.0018 & Small negative \\
Employment (low-poverty counties) & -0.0024 & 0.0010 & 2.017 & -0.0025 & 0.0010 & Null \\
\bottomrule
\end{tabular}
\begin{tablenotes}[flushleft]
\scriptsize
\item \textit{Notes:} \textbf{Country:} United States. \textbf{Research question:} Does a permanent increase in SNAP benefits reduce employment in low-wage industries and reallocate workers toward higher-paying sectors in high-poverty counties? \textbf{Policy mechanism:} The October 2021 Thrifty Food Plan revision permanently raised the maximum SNAP benefit by 21\% (\$36 per person per month), increasing non-labor income for SNAP-eligible households and potentially reducing the opportunity cost of not working in low-wage jobs while enabling job search in better-paying sectors. \textbf{Outcome definition:} Log quarterly beginning-of-quarter employment from the Quarterly Workforce Indicators (QWI/LEHD), measured at the county-industry-quarter level. \textbf{Treatment:} Continuous --- county-level 2019 poverty rate (in percentage points) interacted with a post-October 2021 indicator; higher poverty rates proxy for greater SNAP exposure to the benefit increase. \textbf{Data:} Census QWI (LEHD) county $\times$ NAICS sector $\times$ quarter panel, 2018Q1--2023Q4, merged with 2019 SAIPE county poverty rates. \textbf{Method:} Continuous-treatment difference-in-differences with county $\times$ industry and industry $\times$ quarter fixed effects; standard errors clustered at the state level. \textbf{Sample:} US counties with non-suppressed QWI employment in 7 NAICS sectors (food services, retail, healthcare, admin/waste, manufacturing, professional services, finance); ages 25--54; balanced panel subset used for robustness. SDE $= \hat{\beta} \times \text{SD}(X) / \text{SD}(Y)$ where SD($X$) is the standard deviation of county poverty rate and SD($Y$) is the pre-treatment standard deviation of log employment. Classification refers to magnitude, not statistical significance: Large ($|$SDE$| > 0.15$), Moderate ($0.05$--$0.15$), Small ($0.005$--$0.05$), Null ($< 0.005$).
\end{tablenotes}
\end{threeparttable}
\end{adjustbox}
\end{table}
